\chapter{Autovalores e Autovetores}
\label{cap:autovalores-e-autovetores}

\textbf{Importante:} este capítulo é baseado em \cite{slide7doyuri}.

Neste capítulo, nós começaremos a estudar dois dos conceitos mais importantes da Álgebra Linear: autovalores e autovetores de uma matriz. O seu estudo, historicamente, iniciou-se com o estudo de Equações Diferenciais Ordinárias (EDOs), e possui, além destas, várias outras aplicações, como o estudo de sequências definidas recursivamente, modelagem de epidemias e muitas outras coisas.

\section{Definição de Autovalor e Autovetor}

Nós começaremos esta seção definindo formalmente os conceitos de autovalor e autovetor de uma matriz.

\begin{Definition}[Autovalor e Autovetor]
	Seja $A$ uma matriz quadrada $n\times n$ qualquer. Dizemos que $\lambda\in\mathbb{C}$ é um \textbf{autovalor} de $A$ se existe um vetor $x\in\mathbb{C}^n$ não-nulo tal que:
	
	$$
	Ax = \lambda x.
	$$
	Neste caso, o vetor $x$ é dito ser um \textbf{autovetor} de $A$ associado ao autovalor $\lambda$.
\end{Definition}

Uma das primeiras coisas que decorre desta definição é a seguinte proposição:

\begin{Proposition}
	Seja $A$ uma matriz quadrada $n\times n$ tal que $x\in\mathbb{C}^n$ é um autovetor de $A$. Então, para qualquer escalar $\alpha$, $\alpha x$ também é autovetor de $A$. 
\end{Proposition}

\begin{Proof}
	Se $x$ é um autovetor de $A$, então existe um escalar $\lambda$ tal que $Ax = \lambda x$. Por sua vez, como $A$ é linear, nós temos que:
	
	$$
	A(\alpha x) = \alpha(Ax) = \alpha(\lambda x) = \lambda(\alpha x).
	$$
	
	\qed
\end{Proof}

Uma outra coisa que poderíamos pensar imediatamente é: se uma matriz $A$ possui autovalores iguais a zero, existe alguma relação entre os seus autovetores associados e o núcleo de $A$? A resposta é sim, como veremos agora:

\begin{Proposition}
	Seja $A$ uma matriz $n\times n$. Se $0$ é um autovalor de $A$, então os seus autovetores associados pertencem ao núcleo de $A$. Além disso, $A$ é singular se, e somente se, $0$ é um autovalor de $A$.
\end{Proposition}

\begin{Proof}
	Suponha primeiramente que $0$ é um autovalor de $A$. Então, existe pelo menos um vetor $x$ não-nulo tal que $Ax = 0x = 0$. Ou seja, $x\in N(A)$. Consequentemente, $A$ é singular. Reciprocamente, se $A$ é singular, então existe pelo menos um vetor $x$ não-nulo tal que $x\in N(A)$. Ou seja, $Ax = 0$, o que implica que $Ax = 0x$, ou seja, $0$ é um autovalor de $A$.
	
	\qed
\end{Proof}

A partir da definição, nós podemos também estabelecer uma propriedade válida para qualquer matriz de projeção:

\begin{Proposition}
	Seja $P_{n\times n}$ uma matriz de projeção. Então, os seus únicos autovalores são $0$ e $1$.
\end{Proposition}

\begin{Proof}
	Se $P$ é uma matriz de projeção, nós sabemos que $P^2 = P$. Suponha que $\lambda$ é autovalor de $P$. Então, existe $x\in\mathbb{C}^n$ não-nulo tal que:
	
	$$
	Px = \lambda x.
	$$
	Por sua vez, nós temos que:
	
	$$
	P^2x = P(Px) = P(\lambda x) = \lambda (Px) = \lambda^2 x.
	$$
	Como $P$ é uma matriz de projeção e $x$ é autovetor associado a $\lambda$, segue que:
	
	$$
	\lambda^2 x = \lambda x \iff \lambda^2 = \lambda.
	$$
	Segue imediatamente que $\lambda = 0$ ou $\lambda = 1$.
	
	\qed
\end{Proof}

Dado que achamos os autovalores de uma matriz de projeção qualquer, é interessante também encontrar os seus autovetores. Mas é isto é fácil: para $\lambda = 0$, nós temos $Px = 0$. Logo, qualquer vetor diferente de zero em $N(P)$ é autovetor de $P$. Por sua vez, para $x = 1$, nós temos $Px = x$. Logo, qualquer vetor que pertença ao subespaço sobre o qual $P$ está projetando é um autovetor de $P$. Ou, em outras palavras, qualquer vetor que pertença a $C(P)$ é autovetor de $P$.

Antes de mostrarmos como se faz, de maneira geral, o cálculo dos autovalores e dos autovetores de uma matriz, vamos mostrar um exemplo simples:

\begin{Example}
	Seja $A$ a matriz dada por:
	
	$$
	A = \begin{bmatrix}
		0 & 1 \\
		1 & 0
	\end{bmatrix}.
	$$
	É fácil ver que $A^2 = I$. Assim sendo, se $\lambda$ é autovalor de $A$, nós temos que:
	
	$$
	Ax = \lambda x \iff A^2x = x \iff \lambda^2 x = x \iff \lambda^2 = 1 \iff \lambda = \pm 1.
	$$
	Logo, os autovalores de $A$ são $1$ e $-1$. Por sua vez, como $A$ é uma matriz de permutação, nós temos que:
	
	$$
	\begin{bmatrix}
		0 & 1 \\
		1 & 0
	\end{bmatrix}\begin{bmatrix}
		x_1 \\
		x_2
	\end{bmatrix} = \begin{bmatrix}
		x_2 \\
		x_1
	\end{bmatrix}.
	$$
	Para $x = \begin{bmatrix}
		1 \\
		1
	\end{bmatrix}$, nós temos que $Ax = x$. Logo, $\begin{bmatrix}
		1 \\
		1
	\end{bmatrix}$ é um autovetor associado ao autovalor $1$. Por outro lado, para $x = \begin{bmatrix}
		-1 \\
		1
	\end{bmatrix}$, nós temos que $Ax = -x$. Logo, $\begin{bmatrix}
		-1 \\
		1
	\end{bmatrix}$ é um autovetor associado ao autovalor $-1$.
\end{Example}

Repare que os autovetores associados a um autovalor não são únicos. No mesmo exemplo anterior, $\begin{bmatrix}
	2 \\
	2
\end{bmatrix}$ e $\begin{bmatrix}
	-2 \\
	2
\end{bmatrix}$ também são autovetores associados a $1$ e a $-1$, respectivamente. Mas isto é óbvio, dado que nós já provamos este fato. Na próxima seção, por sua vez, nós estudaremos um método para calcular os autovalores e os autovetores de uma matriz $A$, e faremos mais exemplos (em matrizes $2\times 2$, é claro, porque é um sofrimento fazer isto no caso $3\times 3$, pior ainda é fazer isto no caso $4\times 4$, e muito em breve eu vou contar o problema que é o caso $5\times 5$ ou matrizes maiores).

\section{Cálculo de Autovalores e Autovetores}

Nesta seção, nós veremos um método geral para calcular os autovalores e os autovetores de uma matriz $n\times n$ qualquer. Recordando, lembre-se que $\lambda\in\mathbb{C}$ é um autovalor de uma matriz $A_{n\times n}$ se existe um vetor não-nulo $x\in\mathbb{C}^n$ tal que $Ax = \lambda x$. Assim sendo, nós podemos reescrever esta equação como $Ax - \lambda x = 0$, ou seja, $(A - \lambda I)x = 0$. Note que, como $x\neq 0$, $A - \lambda I$ é uma matriz singular. Assim sendo, é válida a seguinte proposição:

\begin{Proposition}
	Seja $A$ uma matriz $n\times n$ qualquer. Então, $\lambda$ é um autovalor de $A$ se, e somente se, $\det(A - \lambda I) = 0$.
\end{Proposition}

Mais geralmente, nós podemos fazer a seguinte definição:

\begin{Definition}[Polinômio característico]
	Seja $A$ uma matriz $n\times n$ qualquer. Então, o polinômio
	
	$$
	p(\lambda) = \det(A - \lambda I)
	$$
	é chamado o \textbf{polinômio característico} de $A$.
\end{Definition}

Encontrar os autovalores de $A$ equivale a encontrar as raízes do polinômio característico de $A$. Se $A$ é uma matriz $n\times n$, este polinômio possui $n$ raízes (não necessariamente todas distintas) e, portanto, a matriz $A$ possui $n$ autovalores (que podem ser reais ou complexos, como já vimos, e podem ser repetidos).

Na verdade, isto explica o comentário que eu fiz no final da seção anterior: existe um teorema de Álgebra que diz que não existe uma fórmula algébrica fechada para a resolução de equações polinomiais de grau maior ou igual a $5$. No entanto, encontrar os autovalores de uma matriz $A_{n\times n}$ com $n\geq 5$ equivale a resolver uma equação polinomial de grau $n$. Assim sendo, todo e qualquer método para encontrar os autovalores de uma matriz é, necessariamente, um método numérico. Mas não se preocupem com isto agora, nós trabalharemos apenas com o caso $2\times 2$, que é fácil de lidar.

Uma vez que nós encontramos os autovalores de $A$, basta resolver o sistema linear $(A - \lambda I)x = 0$ usando eliminação para encontrar os seus autovetores. Note que nós queremos encontrar algum vetor $x\in\mathbb{C}^n$ tal que $x\in N(A - \lambda I)$. A partir disso, podemos fazer uma definição:

\begin{Definition}[Autoespaço]
	Seja $A$ uma matriz quadrada $n\times n$ qualquer. O subespaço $N(A - \lambda I)$ é chamado o \textbf{autoespaço associado ao autovalor $\lambda$}, e é denotado por $S_{\lambda}$.
\end{Definition}

Vamos agora fazer alguns exemplos para fixar melhor as ideias que nós desenvolvemos até aqui:

\begin{Example}
	Seja $A = \begin{bmatrix}
		3 & 1 \\
		1 & 3
	\end{bmatrix}$. Vamos encontrar os autovalores de $A$, e autovetores associados a estes autovalores. Para encontrar os autovalores, nós resolveremos a equação
	
	$$
	\det(A - \lambda I) = 0.
	$$
	Nós temos que:
	
	$$
	A - \lambda I = \begin{bmatrix}
		3 - \lambda & 1 \\
		1 & 3 - \lambda
	\end{bmatrix}.
	$$
	Assim sendo:
	
	$$
	p(\lambda) = \det(A - \lambda I) = (3 - \lambda)(3 - \lambda) - 1\cdot 1 = \lambda^2 - 6\lambda + 9 - 1 = \lambda^2 - 6\lambda + 8.
	$$
	Igualando o polinômio característico a zero, nós temos:
	
	$$
	\lambda^2 - 6\lambda + 8 = 0 \iff (\lambda - 2)(\lambda - 4) = 0.
	$$
	Logo, nós temos $\lambda = 2$ ou $\lambda = 4$. Estes são os autovalores de $A$. Para encontrar os autovetores associados, nós resolvemos os sistemas lineares substituindo $\lambda$ pelos valores que encontramos.
	
	Para $\lambda = 2$, nós temos:
	
	$$
	(A - 2I)x = 0 \implies \begin{bmatrix}
		1 & 1 \\
		1 & 1
	\end{bmatrix}\begin{bmatrix}
		x_1 \\
		x_2
	\end{bmatrix} = \begin{bmatrix}
		0 \\
		0
	\end{bmatrix}.
	$$
	Uma possível solução para este sistema é dada pelo vetor $x = \begin{bmatrix}
		1 \\
		-1
	\end{bmatrix}$. Assim sendo, este é um autovetor de $A$ associado a $2$.
	
	Para $\lambda = 4$, nós temos:
	
	$$
	(A - 4I)x = 0 \implies \begin{bmatrix}
		-1 & 1 \\
		1 & -1
	\end{bmatrix}\begin{bmatrix}
		x_1 \\
		x_2
	\end{bmatrix} = \begin{bmatrix}
		0 \\
		0
	\end{bmatrix}.
	$$
	Uma possível solução para este sistema é dada pelo vetor $x = \begin{bmatrix}
		1 \\
		1
	\end{bmatrix}$. Logo, este é um autovetor de $A$ associado a $4$.
	
	Note que, na definição de autoespaço que nós fizemos anteriormente, nós temos:
	
	$$
	S_2 = \operatorname{span}\left(\left\{\begin{bmatrix}
		-1 \\
		1
	\end{bmatrix}\right\}\right)
	$$
	e
	
	$$
	S_4 = \operatorname{span}\left(\left\{\begin{bmatrix}
		1 \\
		1
	\end{bmatrix}\right\}\right).
	$$
\end{Example}

\begin{Example}
	Seja $Q = \begin{bmatrix}
		0 & -1 \\
		1 & 0
	\end{bmatrix}$. Os autovalores de $Q$ são dados pelas soluções da equação
	
	$$
	\det(Q - \lambda I) = 0 \iff (-\lambda)(-\lambda) - 1\cdot(-1) = 0 \iff \lambda^2 + 1 = 0.
	$$
	As soluções desta equação são $\lambda = i$ e $\lambda = -i$ (onde $i$ denota a unidade imaginária). Repare que, agora, os autovalores de $Q$ são números complexos. Mas não fiquem com medo, as contas seguem normalmente. Novamente, nós vamos resolver os sistemas lineares para encontrar autovetores de $Q$.
	
	Para $\lambda = i$, nós temos:
	
	$$
	(Q - iI)x = 0 \implies \begin{bmatrix}
		-i & -1 \\
		1 & -i
	\end{bmatrix}\begin{bmatrix}
		x_1 \\
		x_2
	\end{bmatrix} = \begin{bmatrix}
		0 \\
		0
	\end{bmatrix}.
	$$
	Substituindo $x_1$ por $1$, nós encontramos $x_2 = -i$. Logo, $x = \begin{bmatrix}
		1 \\
		-i
	\end{bmatrix}$ é um autovetor associado a $i$.
	
	Para $\lambda = -i$, nós temos:
	
	$$
	(Q + iI)x = 0 \implies \begin{bmatrix}
		i & -1 \\
		1 & i
	\end{bmatrix}\begin{bmatrix}
		x_1 \\
		x_2
	\end{bmatrix} = \begin{bmatrix}
		0 \\
		0
	\end{bmatrix}.
	$$
	Substituindo $x_1$ por $1$, nós obtemos $x_2 = i$. Logo, $x = \begin{bmatrix}
		1 \\
		i
	\end{bmatrix}$ é um autovetor associado a $-i$.
	
	Note que, com a definição de autoespaço que nós temos, nós temos que, neste caso, 
	
	$$
	S_i = \operatorname{span}\left(\left\{\begin{bmatrix}
		1 \\
		-i
	\end{bmatrix}\right\}\right)
	$$
	e
	
	$$
	S_{-i} = \operatorname{span}\left(\left\{\begin{bmatrix}
		1 \\
		i
	\end{bmatrix}\right\}\right).
	$$
\end{Example}

\begin{Example}
	Seja $B = \begin{bmatrix}
		5 & 2 \\
		0 & 5
	\end{bmatrix}$. Nós temos que os autovalores de $B$ são as soluções de
	
	$$
	\det(B - \lambda I) = 0.
	$$
	Por sua vez, é verdade que
	
	$$
	\det(B - \lambda I) = (5 - \lambda)^2.
	$$
	Logo, nós temos a equação
	
	$$
	(5 - \lambda)^2 = 0 \iff \lambda = 5.
	$$
	Para encontrar um autovetor associado a $5$, nós resolveremos o sistema linear $(A - 5I)x = 0$. Nós temos:
	
	$$
	\begin{bmatrix}
		0 & 2 \\
		0 & 0
	\end{bmatrix}\begin{bmatrix}
		x_1 \\
		x_2
	\end{bmatrix} = \begin{bmatrix}
		0 \\
		0
	\end{bmatrix}.
	$$
	Escolhendo $x_1 = 1$, nós temos que $x_2 = 0$. Logo, o vetor $x = \begin{bmatrix}
		1 \\
		0
	\end{bmatrix}$ é um autovetor de $B$ associado a $5$. Por sua vez, nós temos que, para esta matriz, $S_5 = \operatorname{span}\left(\left\{\begin{bmatrix}
		1 \\
		0
	\end{bmatrix}\right\}\right)$.
\end{Example}

No caso $2\times 2$, existe uma maneira bem fácil de encontrar o polinômio característico de $A$, como veremos na proposição a seguir:

\begin{Proposition}
	Seja $A$ uma matriz $2\times 2$, com
	
	$$
	A = \begin{bmatrix}
		a & b \\
		c & d
	\end{bmatrix}.
	$$
	Então, $p(\lambda) = \lambda^2 - (a + d)\lambda + (ad - bc)$.
\end{Proposition}

\begin{Proof}
	Nós temos que $p(\lambda) = \det(A - \lambda I)$. Por sua vez:
	
	$$
	A - \lambda I = \begin{bmatrix}
		a - \lambda & b \\
		c & d - \lambda
	\end{bmatrix}.
	$$
	Assim sendo, 
	
	\begin{align*}
		\det(A - \lambda I) &= (a - \lambda)(d - \lambda) - bc = \lambda^2 - a\lambda - d\lambda + ad - bc \\ &= \lambda^2 - (a + d)\lambda + (ad - bc).
	\end{align*}
	
	\qed
\end{Proof}

No futuro, nós veremos que o traço de uma matriz quadrada é definido como a soma dos seus elementos na diagonal, e é denotado por $\operatorname{Tr}(A)$. Assim sendo, nós podemos reescrever esta equação que obtivemos como $p(\lambda) = \lambda^2 - \operatorname{Tr}(A)\lambda + \det(A)$.

\section{Propriedades dos Autovalores}

Nesta seção, nós estudaremos algumas propriedades dos autovalores de uma matriz, e também veremos algumas definições e propriedades envolvendo o polinômio característico. Começaremos com a seguinte proposição:

\begin{Proposition}
	Sejam $A$ e $B$ duas matrizes $n\times n$ tais que $Ax = \lambda x$ e $Bx = \mu x$, para algum $x\in\mathbb{C}^n$ não-nulo. Então, $\lambda + \mu$ é autovalor de $A + B$. 
\end{Proposition}

\begin{Proof}
	Se $A$ e $B$ satisfazem a propriedade do enunciado, então: 
	
	$$
	Ax + Bx = \lambda x + \mu x \implies (A + B)x = (\lambda + \mu)x.
	$$
	Logo, $\lambda + \mu$ é autovalor de $A + B$.
	
	\qed
\end{Proof}

Por sua vez, em geral, o autovalor da soma de duas matrizes não é a soma dos seus autovalores. Note que, no enunciado anterior, as duas matrizes possuem o mesmo autovetor, o que garante que o autovalor da soma das matrizes seja a soma dos seus autovalores. De igual modo, o autovalor do produto entre duas matrizes não necessariamente é o produto dos seus autovalores.

Agora, veremos o que acontece com os autovalores das potências de uma matriz:

\begin{Proposition}
	Seja $A$ uma matriz $n\times n$ tal que $Ax = \lambda x$, com $x\in\mathbb{C}^n$ não-nulo. Então, $A^{k}x = \lambda^{k}x$.
\end{Proposition}

\begin{Proof}
	Nós provaremos esta proposição por indução em $k$. Para $k = 1$, o resultado é trivial. Para $k = 2$, nós temos que $A^2x = A(Ax) = A(\lambda x) = \lambda(Ax) = \lambda^2 x$. Suponha que a propriedade vale para $k = n - 1$. Provaremos que vale também para $k = n$. De fato, $A^{n}x = A(A^{n - 1}x) = A(\lambda^{n - 1}x) = \lambda^{n - 1}(Ax) = \lambda^n x$.
	
	\qed
\end{Proof}

A próxima propriedade relaciona os autovalores de uma matriz não-singular com os autovalores da sua inversa:

\begin{Proposition}
	Seja $A$ uma matriz $n\times n$ invertível tal que $Ax = \lambda x$, para algum $x\in\mathbb{C}^n$ não-nulo. Então, $A^{-1}x = \dfrac{1}{\lambda}x$.
\end{Proposition}

\begin{Proof}
	Primeiramente, como $A$ é invertível, nós temos que $x$ pode ser escrito como:
	
	$$
	x = A^{-1}Ax.
	$$
	Como $x$ é autovetor de $A$, segue que:
	
	$$
	x = A^{-1}(Ax) = A^{-1}(\lambda x) = \lambda A^{-1}x.
	$$
	Segue que
	
	$$
	A^{-1}x = \dfrac{1}{\lambda}x.
	$$
	
	\qed
\end{Proof}

Agora, nós vamos tratar do polinômio característico de uma matriz $A_{n\times n}$. Recordando a sua definição, nós temos que o polinômio característico de $A$ é definido como $p(\lambda) = \det(A - \lambda I)$. Por sua vez, da linearidade do determinante, nós podemos ver facilmente que $\det(k\cdot A) = k^n\det(A)$, para algum escalar $k$. Assim sendo, uma outra forma de enxergar o polinômio característico é como $p(\lambda) = (-1)^n\det(\lambda I - A)$. Note que, nesta segunda maneira de ver o polinômio, as suas raízes não se alteram. Como este é um polinômio de grau $n$, nós também podemos escrevê-lo, nesta segunda forma, como:

$$
p(\lambda) = (-1)^n (\lambda^n + c_{n - 1}\lambda^{n - 1} + \cdots + c_1\lambda + c_0).
$$

Por sua vez, se $\lambda_1, \cdots, \lambda_n$ são as raízes de $p$ (incluindo as multiplicidades), pelo Teorema Fundamental da Álgebra (que não enunciaremos nem provaremos aqui), nós temos que $p$ pode ser decomposto como:

$$
p(\lambda) = (-1)^n(\lambda - \lambda_1)(\lambda - \lambda_2)\cdots(\lambda - \lambda_n).
$$

Usando esta decomposição, nós poderemos provar duas propriedades. A primeira delas é a seguinte:

\begin{Proposition}
	Sejam $\lambda_1, \cdots, \lambda_n$ os autovalores de $A$. Então, $\det(A) = \prod\limits_{i = 1}^{n} \lambda_i$.
\end{Proposition}

\begin{Proof}
	Nós temos que $\det(A) = \det(A - 0I) = p(0)$. Assim sendo, vamos calcular $p(0)$:
	
	$$
	p(0) = (-1)^n\prod\limits_{i = 1}^{n} (0 - \lambda_i) = (-1)^n\prod\limits_{i = 1}^{n} (-\lambda_i) = (-1)^n(-1)^n\prod\limits_{i = 1}^{n} \lambda_i = (-1)^{2n}\prod\limits_{i = 1}^{n} \lambda_i = \prod\limits_{i = 1}^{n} \lambda_i.
	$$
	
	\qed
\end{Proof}

A partir deste fato, nós podemos determinar facilmente qual é o coeficiente $c_0$ do polinômio expandido. De fato, nós temos que:

$$
p(0) = (-1)^n(0^{n} + c_{n - 1}\cdot0^{n - 1} + \cdots + c_1\cdot 0 + c_0).
$$
Assim sendo, $p(0) = (-1)^n c_0$ e, como $p(0) = \det(A)$, segue que $c_0 = (-1)^n\det(A)$.

Antes de demonstrarmos a segunda propriedade, vamos fazer uma definição:

\begin{Definition}[Traço de uma matriz]
	Seja $A$ uma matriz $n\times n$ tal que a sua entrada $ij$ é dada por $a_{ij}$. Então, definimos o \textbf{traço} de $A$, com notação $\operatorname{Tr}(A)$, por:
	
	$$
	Tr(A) = \sum\limits_{i = 1}^{n} a_{ii}.
	$$
\end{Definition}

Em resumo, o traço de uma matriz $A$ é a soma dos seus elementos diagonais. Agora, nós podemos enunciar e demonstrar a segunda propriedade:

\begin{Proposition}
	Seja $A$ uma matriz $n\times n$ com autovalores $\lambda_1, \cdots, \lambda_n$. Então, $Tr(A) = \sum\limits_{i = 1}^{n}$.
\end{Proposition}

\begin{Proof}
	Se $p(\lambda)$ é o polinômio característico de $A$, então nós calcularemos o coeficiente $c_{n - 1}$ de duas formas. Primeiramente, nós podemos enxergar $p(\lambda)$ como:
	
	$$
	p(\lambda) = (-1)^n(\lambda - \lambda_1)\cdots(\lambda - \lambda_n).
	$$
	O coeficiente $c_{n - 1}$ é aquele que acompanha o termo $\lambda^{n - 1}$. Assim sendo, como $p(\lambda)$ possui $n$ termos ``entre parênteses'', nós precisamos ``pegar'' $n - 1$ $\lambda$'s de parênteses diferentes, e escolher $1$ termo diferente de $\lambda$ em um dos parênteses. Por sua vez, como nós precisamos fazer isto de todas as formas possíveis, ao final nós chegamos a $c_{n - 1} = - (\lambda_1 + \cdots + \lambda_n)$.
	
	Por outro lado, nós podemos enxergar $p(\lambda)$ como $(-1)^n\det(\lambda I - A)$, e analisar o que ocorre com $\det(\lambda I - A)$. Nós temos que a expansão por cofatores nos dá:
	
	\begin{align*}
		\det(\lambda I - A) &= \det\left(\begin{bmatrix}
			\lambda - a_{11} & a_{12} & \cdots & a_{1n} \\
			a_{21} & \lambda - a_{22} & \cdots & a_{2n} \\
			\vdots & \vdots & \ddots & \vdots \\
			a_{n1} & a_{n2} & \cdots & \lambda - a_{nn}
		\end{bmatrix}\right)\\ &= (\lambda - a_{11})\det\left(\begin{bmatrix}
			\lambda - a_{22} & \cdots & a_{2n} \\
			\vdots & \ddots & \vdots \\
			a_{n2} & \cdots & \lambda - a_{nn}
		\end{bmatrix}\right) + q(\lambda),
	\end{align*}
	onde $q$ é um polinômio de grau $n - 2$. Por sua vez, fazendo mais uma expansão, nós temos que:
	
	\begin{align*}
		& (\lambda - a_{11})\det\left(\begin{bmatrix}
			\lambda - a_{22} & \cdots & a_{2n} \\
			\vdots & \ddots & \vdots \\
			a_{n2} & \cdots & \lambda - a_{nn}
		\end{bmatrix}\right) + q(\lambda)\\ = & (\lambda - a_{11})(\lambda - a_{22})\det\left(\begin{bmatrix}
			\lambda - a_{33} & \cdots & a_{3n} \\
			\vdots & \ddots & \vdots \\
			a_{n3} & \cdots & \lambda - a_{nn}
		\end{bmatrix}\right) + \tilde{q}(\lambda),
	\end{align*}
	onde $\tilde{q}$ é um outro polinômio de grau $n - 2$. Executando este raciocínio recursivamente, ao final nós obtemos:
	
	$$
	\det(\lambda I - A) = (\lambda - a_{11})(\lambda - a_{22})\cdots(\lambda - a_{nn}) + \hat{q}(\lambda),
	$$
	onde $\hat{q}$ é um polinômio de grau $n - 2$. Usando o mesmo raciocínio que nós usamos no primeiro caso, nós temos que o coeficiente $c_{n - 1}$ é dado por $c_{n - 1} = - (a_{11} + a_{22} + \cdots + a_{nn})$. Logo, segue que:
	
	$$
	a_{11} + a_{22} + \cdots + a_{nn} = \lambda_1 + \cdots + \lambda_n.
	$$
	Mas $a_{11} + a_{22} + \cdots + a_{nn} = \operatorname{Tr}(A)$. Logo, $\operatorname{Tr}(A) = \lambda_1 + \cdots + \lambda_n$.
	
	\qed
\end{Proof}